% =====================================================================
\section{FTN-LoRA: $B_0$ initialization, ablations, and the routing-vs-random control}\label{app:lora_ablations}

This appendix documents two design choices for the LoRA-bank experiment in \cref{sec:lora}: (i)~the non-standard initialization $B_0\!\sim\!\mathcal{N}(0,0.01)$ in place of the usual $B_0{=}0$, and (ii)~the random-per-batch-mask ablation that isolates the contribution of the support-conditioned configurer from the rest of the FTN-LoRA training pipeline. \Cref{tab:lora_full} reports the full ablation, including both diagnostic $B_0{=}0$ runs and the headline five-seed results.

\paragraph{Standard $B_0{=}0$ initialization.}
Standard LoRA initializes $A$ from a Kaiming-uniform distribution and $B$ to all zeros, so the adapter contributes nothing on step~0 and the fine-tuned model starts at the base backbone exactly. We first ran FTN-LoRA with this initialization. Retention was already reasonable (FTN-LoRA $B_0{=}0$ + Router reaches ACC $0.76\pm0.04$, FM $0.19\pm0.02$ across two seeds), but the configurer's head-start NLL was substantially worse than the retrieval baselines: $1.66\pm0.17$ (FTN $B_0{=}0$) and $1.95\pm0.35$ (FTN $B_0{=}0$ + Router) versus $1.22$--$1.32$ for the matched-backbone baselines. The configurer was producing masks that actively hurt the new task.

\paragraph{Diagnosis.}
With $B_j$ exactly zero, atom $j$'s contribution to the bank output is $x A_j B_j = 0$ irrespective of $A_j$, so $\partial \mathcal{L}/\partial m_j = 0$. Under the AnnealSTE straight-through estimator, the gradient on the mask logit for atom $j$ is therefore also zero: the inner loop receives no signal about whether atom $j$ would help on the new task. It can deselect atoms whose previous-task training is actively misaligned with the support set, but it cannot pull useful untrained atoms into the active set without first selecting them by chance. The parallel-MLP backbone in \cref{sec:experiments} avoids this because every parallel MLP, trained or not, contributes a stochastic non-zero output, so the configurer's gradient on the mask is informative about every atom from step~0.

\paragraph{$B_0\!\sim\!\mathcal{N}(0,0.01)$.}
We replace zero initialization with a small Gaussian, $B_0\!\sim\!\mathcal{N}(0,0.01)$. The model still starts close to the base backbone (the per-atom contribution is small), but every atom produces a non-zero output that gives the configurer a non-zero gradient signal on its mask logit. This change unlocks the FTN-LoRA family results in \cref{tab:lora_full}: ACC rises from $0.62$ to $0.83$, FM falls from $0.28$ to $0.09$, and BWT flips from $-0.20$ to $+0.19$, all within multi-seed sample noise of the random-per-batch-mask ablation that shares the same training pipeline (next paragraph).

\paragraph{The random-per-batch-mask ablation.}
A reasonable concern is that the small non-zero $B_0$ injects useful adapter perturbations on its own, and that the headline FTN-LoRA gains in \cref{sec:lora} therefore reflect a non-FTN cause. We test this with a five-seed control: \emph{FTN-LoRA, random per-batch masks} keeps the FTN-LoRA training pipeline (mask-overlap-filtered replay, $B_0\!\sim\!\mathcal{N}(0,0.01)$, the same $100$-step schedule and active-parameter budget, the same per-batch reconfiguration cadence) but replaces the support-conditioned configurer with a fresh random $k$-of-$H$ mask each step. The eval-time configurer is unchanged. The result decomposes the headline cleanly: the random-mask ablation matches FTN-LoRA on retention (ACC $0.85\pm0.12$ vs $0.83\pm0.10$, FM $0.12\pm0.07$ vs $0.09\pm0.09$) and on head-start NLL ($1.38\pm0.18$ vs $1.42\pm0.25$), so the training pipeline rather than the configurer is the source of those numbers; but it does not flip backward transfer ($-0.08\pm0.12$ vs $+0.19\pm0.10$), so the support-conditioned configurer specifically is what produces FTN-LoRA's positive BWT. The FTN-LoRA mean lies more than one standard deviation above zero across five seeds and the random-mask mean lies below zero, so the sign of BWT separates the two configurations on this seed budget; we do not claim a paired-bootstrap significance test here, and the per-method standard deviations should be read as spread rather than as a tight confidence interval.

\paragraph{What this jointly establishes and what it does not.}
The two ablations together support a narrower mechanism claim than the controlled MLP suite's: at LoRA scale the support-conditioned configurer is the load-bearing piece for backward transfer specifically, while the headline retention is explained by the training pipeline (filtered replay, non-zero $B_0$, per-batch reconfiguration on a wide bank). The LoRA experiment is silent on head-start: every method in \cref{tab:lora_full} lands in the narrow band $1.22$--$1.42$, well within the per-row standard deviations of $0.15$--$0.25$. Head-start can only vary across methods if prior-task adapters carry usable bias for the new task, and the six-task stream was selected for diversity (Chinese stance, English finance sentiment, English question type, English emotion, English news topic, English ontology) rather than for measured cross-task structure: most pairings cross both a language change and a label-space change, so even an oracle would have limited prior-update structure to extract. The parallel-MLP suite avoids this confound by construction (correlated synthetic and shared-shuffle MNIST have hand-tuned concept overlap), and remains the appropriate probe of FTN's head-start claim. BWT separates on the LoRA suite because it depends on training-time bank dynamics across tasks rather than on any specific task pair's cross-task structure. What we have not done at LoRA scale: a $B_0$ scale sweep, a Qwen-7B/14B scale-up, an order-sensitivity check, or an LLM stream with measurable cross-task structure that would let head-start adjudicate the configurer's pre-update mechanism.

\begin{table}[t]
\centering
\small
\caption{Full FTN-LoRA ablation table. Headline rows (SeqLoRA, SeqLoRA + ER, Indep.\ LoRA, Indep.\ LoRA + Router, FTN-LoRA random per-batch masks, FTN-LoRA) are computed over five seeds; the three diagnostic FTN rows that use standard LoRA initialisation ($B_0{=}0$) are computed over two seeds and are reported only as the diagnostic that motivated the $B_0{\sim}\mathcal{N}(0,0.01)$ choice. RandMask + Router is two seeds. The diagnostic narrative is in \cref{app:lora_ablations}.}
\label{tab:lora_full}
\begin{tabular}{lcccc}
\toprule
Method & ACC $\uparrow$ & FM $\downarrow$ & BWT $\uparrow$ & HS NLL $\downarrow$ \\
\midrule
SeqLoRA & 0.417\,\scriptsize$\pm$0.131 & 0.550\,\scriptsize$\pm$0.197 & -0.465\,\scriptsize$\pm$0.174 & 1.37\,\scriptsize$\pm$0.16 \\
SeqLoRA + ER & 0.750\,\scriptsize$\pm$0.132 & 0.145\,\scriptsize$\pm$0.112 & -0.090\,\scriptsize$\pm$0.129 & 1.37\,\scriptsize$\pm$0.22 \\
RandMask + Router & 0.573\,\scriptsize$\pm$0.192 & 0.287\,\scriptsize$\pm$0.053 & -0.263\,\scriptsize$\pm$0.088 & 1.25\,\scriptsize$\pm$0.13 \\
Indep.\ LoRA (oracle) & 0.392\,\scriptsize$\pm$0.139 & 0.275\,\scriptsize$\pm$0.158 & -0.255\,\scriptsize$\pm$0.174 & 1.22 \\
Indep.\ LoRA + Router & 0.642\,\scriptsize$\pm$0.159 & 0.150\,\scriptsize$\pm$0.126 & -0.105\,\scriptsize$\pm$0.112 & 1.26\,\scriptsize$\pm$0.17 \\
FTN-LoRA, $B_0=0$, no replay & 0.729 & 0.225\,\scriptsize$\pm$0.035 & -0.087\,\scriptsize$\pm$0.124 & 3.49\,\scriptsize$\pm$2.72 \\
FTN-LoRA, $B_0=0$ + Router & 0.760\,\scriptsize$\pm$0.044 & 0.188\,\scriptsize$\pm$0.018 & -0.125 & 1.95\,\scriptsize$\pm$0.35 \\
FTN-LoRA, $B_0=0$ & 0.615\,\scriptsize$\pm$0.162 & 0.275\,\scriptsize$\pm$0.106 & -0.200\,\scriptsize$\pm$0.106 & 1.66\,\scriptsize$\pm$0.17 \\
FTN-LoRA, random per-batch masks & 0.854\,\scriptsize$\pm$0.118 & 0.120\,\scriptsize$\pm$0.074 & -0.075\,\scriptsize$\pm$0.115 & 1.38\,\scriptsize$\pm$0.18 \\
FTN-LoRA & 0.833\,\scriptsize$\pm$0.101 & 0.085\,\scriptsize$\pm$0.086 & 0.190\,\scriptsize$\pm$0.101 & 1.42\,\scriptsize$\pm$0.25 \\
\bottomrule
\end{tabular}
\end{table}


\begin{table}[t]
\centering
\small
\caption{Per-task head-start NLL on the six-task multilingual classification stream described in \cref{sec:lora}. Each entry is the answer-token NLL on the held-out batch of that task, evaluated under the method's mask BEFORE any gradient update on the task. Lower is better. Tasks are ordered as canonically presented during training; task~1 (C-STANCE-A) is omitted because no prior context exists.}
\label{tab:lora_hs_per_task}
\begin{tabular}{lccccc}
\toprule
Method & financial & trec & emotion & ag news & dbpedia \\
\midrule
SeqLoRA & 2.36 & 0.60 & 1.06 & 0.21 & 2.60 \\
SeqLoRA + ER & 2.60 & 0.56 & 1.11 & 0.22 & 2.39 \\
RandMask + Router & 1.96 & 0.35 & 0.81 & 0.21 & 2.93 \\
Indep.\ LoRA (oracle) & 1.47 & 0.48 & 1.84 & 0.16 & 2.17 \\
Indep.\ LoRA + Router & 2.36 & 0.42 & 1.22 & 0.13 & 2.15 \\
FTN-LoRA, $B_0=0$, no replay & 1.61 & 9.29 & 1.77 & 0.45 & 4.33 \\
FTN-LoRA, $B_0=0$ + Router & 1.99 & 2.87 & 1.87 & 0.36 & 2.66 \\
FTN-LoRA, $B_0=0$ & 2.24 & 0.94 & 1.13 & 0.38 & 3.61 \\
FTN-LoRA, random per-batch masks & 1.41 & 0.75 & 1.08 & 0.33 & 3.35 \\
FTN-LoRA & 1.81 & 0.67 & 1.18 & 0.37 & 3.08 \\
\bottomrule
\end{tabular}
\end{table}

